An observer looks at a periodic ripple through a lens of finite resolution $\sigma$ and is free to choose where and how sharply to focus.
We prove that, for every window whose ripple is weak in the sup norm ($\eta\le1/\ell$, $\ell=\ln(1/\sigma)$), no choice of focus beats the logarithmic law of the resolution gap:
wherever the observer focuses, and however fine the resolution, $g\sqrt{\ell}\ge1-\bar\varepsilon$ for every state, for all $\ell\ge55.26$ (that is, $\sigma\le e^{-55.26}\approx1.002\times10^{-24}$, with no lower limit on $\sigma$); for $\ell\ge960.8$ the bound is $(1-\bar\varepsilon)\sqrt{1+10/\ell}$.
Here $\bar\varepsilon$ bounds the remainder on the window class, as in Papers~F and~H~\cite{F,H}.
Paper~H proved this for the states whose focus spans at least about $0.7$ ripple periods and for $\ell\le960.8$.
The present paper closes the rest of the map with four theorems, one for each kind of focus.
For a focus at the resolution itself, positivity alone leaves the two background terms intact, and $g$ stays of order one \proofst.
For a narrow focus, and for a wide focus, a local expansion with explicit remainders and one brightness hypothesis give the bound; the constants are certified by interval arithmetic \status{proof, computer-assisted}.
For a very wide focus, which averages hundreds of ripple periods, two elementary facts suffice \proofst:
the cross term of the curvature step, measured in mean square over the ripple rather than at its worst point, costs eight percent of the ripple information whatever the width of the focus;
and a weight hidden at a zero of the phase mismatch reappears, by periodicity, one and two periods away, where the mismatch does not vanish.
This part uses no bound on $\eta$, so the very wide focus is settled for every window of the class.
We state exactly what is proved, what is computer-assisted, and what is only numerical, and we list what remains outside: windows with a stronger ripple, the class constant $\bar\varepsilon$, and the upper side.
